% Generated draft: requires mathematical and source review.
\documentclass[11pt]{article}
\usepackage[margin=25mm]{geometry}
\usepackage{fontspec}
\setmainfont{Latin Modern Roman}
\usepackage{amsmath,amssymb,mathtools,unicode-math}
\setmathfont{Latin Modern Math}
\usepackage{enumitem,needspace,xurl,hyperref}
\setlength{\parindent}{0pt}
\setlength{\parskip}{5pt}
\setlength{\emergencystretch}{3em}
\newcommand{\N}{\mathbb{N}}
\newcommand{\Z}{\mathbb{Z}}
\newcommand{\Q}{\mathbb{Q}}
\newcommand{\R}{\mathbb{R}}
\newcommand{\C}{\mathbb{C}}
\newcommand{\F}{\mathbb{F}}
\newcommand{\A}{\mathbb{A}}
\newcommand{\PP}{\mathbb{P}}
\newcommand{\E}{\mathbb{E}}
\newcommand{\eps}{\varepsilon}
\newcommand{\dd}{\,\mathrm{d}}
\begin{document}
\section*{UnsolvedMath ID: 30000841}
\textbf{Source problem code:} catalogue entry
\subsection{Definitions and mathematical statement}
Let \((M,\mathcal F)\) be a Riemannian foliation in the usual sense that its normal bundle has a holonomy-invariant transverse metric. Choose the leafwise metric and analytic setting intended in the source. For a leaf \(L\) and form degree \(q\), let
\[
\Delta_{L,q}=d_Ld_L^*+d_L^*d_L,
\qquad
\mathcal H^q_{(2)}(L)=\ker\Delta_{L,q}\subseteq L^2\Omega^q(L),
\]
where an appropriate self-adjoint realization is required, and let \(P_{L,q}\) be the orthogonal projection onto this harmonic space. Its Schwartz kernel represents the projection as an integral operator on leafwise forms when the relevant kernel interpretation is available.

The question in both the original and cleaned statements is: does the family of projections onto leafwise harmonic forms always have a Schwartz kernel smooth in the transverse directions? In a local realization of the kernel family with transverse parameter \(t\), the intended property is smooth dependence on \(t\), in addition to the leafwise regularity. The exact transverse-smoothness topology, groupoid or covering space carrying the kernels, form-degree quantification, and analytic realization are not specified in this extraction.

The Riemannian-foliation assumption is retained. It is not replaced by the much weaker assertion that each leaf is given a Riemannian metric. This distinction is material: the supplied literature assessment says transverse smoothness is already proved for Riemannian foliations and certain compact-leaf or Hausdorff-groupoid settings, while calling the case of arbitrary foliations open. That assessment concerns a broader question than the literal extracted statement and conflicts with treating the latter as an unqualified open problem. The intended target must therefore be checked rather than silently broadened.
\subsection{Short English statement}
Must the leafwise harmonic projection of a Riemannian foliation have a transversely smooth Schwartz kernel? The supplied assessment says this Riemannian case is proved and instead discusses an open question for arbitrary foliations, so the intended scope needs review.
\subsection{Sources}
\begin{itemize}
\item Local UnsolvedMath registry, permanent ID \(30000841\), source code catalogue entry, ``Transverse Smoothness of Leafwise Harmonic Projections.'' No individual proposer is supplied.
\item \emph{Noncommutative Geometry} (2008), Oberwolfach Reports, DOI 10.4171/owr/2007/43, pp. 2543--2608: \url{https://doi.org/10.4171/owr/2007/43}.
\item Supplied reference, ``twisted higher harmonic signature'': \url{https://arxiv.org/abs/0711.0352}.
\item Supplied literature triage, checked 2026-08-21, labels the item ``open'' but reports the Riemannian-foliation and several compact-leaf or Hausdorff-groupoid cases proved, with arbitrary foliations remaining open. This scope discrepancy is unresolved. The dated assessment derives from supplied OWR curation; time-sensitive status and cited claims require independent verification.
\end{itemize}
\end{document}
